\chapter{Introduction}
\label{ch:introduction}

Modern railway systems are undergoing a profound transformation driven by increased traffic volumes, rising safety expectations, and long-term ambitions toward automated and driverless train operation. As highlighted in the RAIL-BENCH kickoff presentation, demographic changes, strong competition from other transportation modes, and the growing demand for both passenger and freight transport intensify the need for reliable perception systems capable of supporting future \acrfull{ato} scenarios. Many of the tasks traditionally performed by human drivers - including obstacle recognition, signal interpretation, and track infrastructure monitoring - are fundamental perception tasks that must be replicated by technical systems to enable safe autonomous operation. Deep learning has demonstrated remarkable progress in visual perception, yet its application in the railway domain remains limited by the scarcity of publicly available datasets. The lack of standardized evaluation protocols also prevents a fair comparison between different methods. \cite{railbench-kickoff}


As explained by Bätz et al.~\cite{railbench-object-paper}, a particular challenge arises from the nature of existing railway infrastructure.
While automated operation has already been successfully established in isolated environments such as metro systems and people movers, these systems operate on infrastructure where access to the tracks is inherently restricted and the track occupancy can be reliably controlled. In contrast, conventional railway networks frequently include level crossings, open platform areas and shared-traffic environments, making it impossible to assume that the track is free of obstacles at all times. To address these challenges, \acrfull{kit} has introduced \emph{RAIL-BENCH}, the first benchmark suite dedicated to perception tasks in railway environments. RAIL-BENCH provides open datasets, unified evaluation metrics, and clearly defined testing procedures to facilitate fair comparisons between methods and accelerate research and development in this safety-critical domain \cite{railbench-object-paper}. RAIL-BENCH supports five scientific challenges \cite{railbench-object-paper}:
\begin{itemize}
    \item Rail Track Detection
    \item Object Detection
    \item Vegetation Segmentation
    \item Multiple Object Tracking
    \item Monocular Visual Odometry
\end{itemize}

Among the five challenges included in the suite, the \emph{RAIL-BENCH Object Detection Challenge} focuses on detecting and classifying objects relevant to railway operation. This challenge is built upon a dataset of 2500 real-world RGB images recorded from a train driver's perspective under diverse environmental conditions, including urban, rural, and station scenarios, and varying weather and camera setups \cite{railbench-object-detection-page}. Each image is annotated with axis-aligned bounding boxes covering seven object categories \cite{railbench-object-detection-page}:
\begin{itemize}
    \item Train
    \item Person
    \item Bicycle
    \item Road Vehicle
    \item Signal
    \item Signal Pole
    \item Catenary Pole
\end{itemize}
Additional annotation attributes such as occlusion levels and ignore flags further reflect the complexity of real railway scenes, where small, distant, or densely grouped objects frequently occur.

The challenge evaluates submitted predictions using the COCO-style \acrfull{map} across multiple \acrfull{iou} thresholds. Three difficulty levels - easy, moderate, and hard - are defined based on bounding box size and occlusion, with final rankings determined by the \acrshort{map} score in the hard setting averaged over all \acrshort{iou} thresholds and object classes \cite{railbench-object-detection-page}. This evaluation protocol ensures that models are assessed not only on clearly visible objects but also on small and heavily occluded instances.
\\

The goal of our project is to develop, train, and evaluate an object detection model that meets the requirements of the RAIL-BENCH Object Detection Challenge. Working with the official dataset and adhering strictly to the challenge guidelines, we aim to investigate the performance of modern deep learning architectures in railway-specific environments, analyze their strengths and limitations, and contribute to the broader effort of establishing robust perception systems for future autonomous railway operation. This report documents our approach, covering related work, methodology, experimental setup and evaluation results as well as conclusions drawn from our participation in the challenge.
